\documentclass[12pt,a4paper]{article}

% -------------------------------------------------
% Sprache / Kodierung
% -------------------------------------------------
\usepackage[T1]{fontenc}
\usepackage[utf8]{inputenc}
\usepackage[ngerman]{babel}
\usepackage{lmodern}
\usepackage{microtype}

% -------------------------------------------------
% Mathematik
% -------------------------------------------------
\usepackage{amsmath,amssymb,amsthm,mathtools}

% -------------------------------------------------
% Layout / Grafik / Farben
% -------------------------------------------------
\usepackage[a4paper,margin=2.3cm]{geometry}
\usepackage{booktabs}
\usepackage{xcolor}
\usepackage[hidelinks,unicode]{hyperref}
\pdfstringdefDisableCommands{%
  \def\ü{ue}\def\ö{oe}\def\ä{ae}\def\Ü{Ue}\def\Ö{Oe}\def\Ä{Ae}\def\ss{ss}%
}
\usepackage[most]{tcolorbox}
\usepackage{enumitem}
\usepackage{array}
\usepackage{siunitx}

% -------------------------------------------------
% Farben (konsistent mit Vorgängerdokumenten)
% -------------------------------------------------
\definecolor{lückenrot}{RGB}{180,40,40}
\definecolor{bewiesengrün}{RGB}{30,100,50}
\definecolor{warnorange}{RGB}{200,100,10}
\definecolor{hintergrundblau}{RGB}{235,242,250}
\definecolor{hintergrundgelb}{RGB}{255,252,225}
\definecolor{adischviolett}{RGB}{90,20,140}
\definecolor{fehlerrot}{RGB}{200,0,0}
\definecolor{konjviolett}{RGB}{80,0,130}
\definecolor{e8grau}{RGB}{55,55,90}

% -------------------------------------------------
% tcolorbox-Stile
% -------------------------------------------------
\tcbset{
  lücke/.style={colback=red!8, colframe=lückenrot, fonttitle=\bfseries,
    title={Offene Lücke}},
  ergebnis/.style={colback=green!7, colframe=bewiesengrün,
    fonttitle=\bfseries, title={Bewiesenes Ergebnis}},
  warnung/.style={colback=orange!10, colframe=warnorange,
    fonttitle=\bfseries, title={Achtung / Heuristik}},
  adisch/.style={colback=violet!6, colframe=adischviolett,
    fonttitle=\bfseries, title={2-adische Perspektive}},
  kernidee/.style={colback=hintergrundblau, colframe=blue!60!black,
    fonttitle=\bfseries, title={Kernidee}},
  konj/.style={colback=violet!5, colframe=konjviolett, fonttitle=\bfseries,
    title={Konjektur (unbewiesen)}},
  lte/.style={colback=yellow!8, colframe=e8grau, fonttitle=\bfseries,
    title={LTE-Analyse}},
  hauptsatz/.style={colback=green!10, colframe=bewiesengrün!80!black,
    fonttitle=\bfseries\large, title={Stärkster beweisbarer Satz}},
  korrektur/.style={colback=yellow!5, colframe=orange!70!black,
    fonttitle=\bfseries},
  sensitivitaet/.style={colback=blue!5, colframe=blue!50!black,
    fonttitle=\bfseries, title={Sensitivitätsanalyse}},
}

% -------------------------------------------------
% Theorem-Umgebungen
% -------------------------------------------------
\newtheorem{definition}{Definition}[section]
\newtheorem{proposition}[definition]{Proposition}
\newtheorem{theorem}[definition]{Theorem}
\newtheorem{lemma}[definition]{Lemma}
\newtheorem{korollar}[definition]{Korollar}
\newtheorem{bemerkung}[definition]{Bemerkung}
\newtheorem{hypothese}[definition]{Hypothese}
\newtheorem{satz}[definition]{Satz}

\newenvironment{beweis}{\begin{proof}[Beweis]}{\end{proof}}

% -------------------------------------------------
% Makros
% -------------------------------------------------
\newcommand{\vii}{\nu_2}
\newcommand{\U}{\mathcal{U}}
\newcommand{\N}{\mathbb{N}}
\newcommand{\Z}{\mathbb{Z}}
\newcommand{\R}{\mathbb{R}}
\newcommand{\E}{\mathbb{E}}
\newcommand{\Ztwo}{\mathbb{Z}_2}
\newcommand{\abs}[1]{\lvert#1\rvert}
\newcommand{\atwo}[1]{\lvert#1\rvert_2}
\newcommand{\logF}{\log F(B)}

% -------------------------------------------------
% Dokument
% -------------------------------------------------
\title{%
  Konditionierte Drift-Analyse im Lyapunov-Rahmen:\\[0.3em]
  Präzisierung der $\nu$-Verteilung an C-Ketten-Endpunkten\\[0.5em]
  \large Präzisierungs-Appendix zu \textit{collatz\_lyapunov.tex}%
}
\author{Thomas Hoffbauer}
\date{16.\ Juni 2026}

\begin{document}
\maketitle

\begin{abstract}
Dieses Dokument präzisiert die Lyapunov-Drift-Rechnung aus
\textit{collatz\_lyapunov.tex}.
Der naive Ansatz $\E[\nu] \approx 3$ (geometrische Verteilung) ist für
C-Ketten-Endpunkte zu optimistisch: Das Lifting-the-Exponent-Lemma (LTE)
zeigt, dass $\vii(3n+1)$ an solchen Endpunkten typischerweise
$O(\log l)$ statt $\approx 3$ beträgt.
Die korrekte Driftrechnung muss daher \emph{konditioniert} nach der
C-Kettenlänge $l$ aufgestellt werden.

Das Hauptergebnis ist heuristisch: Auch nach Konditionierung ergibt sich
$\E[\logF] \approx -0{,}33 < 0$, der Drift bleibt negativ.
Eine Sensitivitätsanalyse zeigt, dass dieser Schluss für
$\E[\nu \mid l{=}k] < 3{,}34$ robust ist.
Abschließend werden die verbleibenden rigorosen Lücken klar benannt.

\medskip
\noindent\textbf{Status:}
\textcolor{bewiesengrün}{Schritt~1--2 (Wahrscheinlichkeiten, LTE-Schranke): bewiesen.}
\textcolor{warnorange}{Schritte~3--5 (konditionierte Driftrechnung): heuristisch.}
\textcolor{lückenrot}{Schritt~6 (Ergodizität, Zyklenfreiheit): offen.}
\end{abstract}

\tableofcontents
\bigskip

\begin{tcolorbox}[warnung, title={Einordnung dieses Dokuments}]
Die Collatz-Vermutung ist unbewiesen.
Dieser Text trennt konsequent zwischen
\textbf{bewiesenen Aussagen} (grün),
\textbf{Heuristiken} (orange)
und \textbf{offenen Lücken} (rot).
Die konditionierte Drift-Rechnung ist eine \emph{heuristische Verbesserung},
kein Beweis der Collatz-Vermutung.
\end{tcolorbox}

% ====================================================
\section{Motivation: Warum die naive Rechnung zu optimistisch ist}
\label{sec:motivation}
% ====================================================

In der Lyapunov-Analyse wird der logarithmische Drift eines
EABC-Blocks~$B$ betrachtet.
Ein Block der Kettenlänge $l \geq 0$ führt eine C-Kette der Länge $l$
aus, gefolgt von einem EA-Schritt.
Die zentrale Größe ist der \emph{Block-Faktor}
\[
  F(B) = \frac{3^{l+1}}{2^{l+\nu}},
\]
wobei $\nu = \vii(3n_l + 1)$ die 2-adische Bewertung am Block-Endpunkt ist.

\begin{tcolorbox}[lte]
\textbf{Naive Annahme (zu optimistisch):}
Ohne Beachtung der Block-Struktur gilt für beliebige ungerade~$n$:
$\vii(3n+1) \geq 1$ mit $P(\vii(3n+1) = k) = 2^{-(k-1)}$ für $k \geq 1$,
also $\E[\nu] = 3$.

\medskip
\textbf{Das Problem:}
Am Block-Endpunkt einer C-Kette der Länge $l \geq 1$ hat~$n_l$ die Form
\[
  n_l = 2 \cdot 3^l \cdot m - 1
\]
für ein ungerades $m$.
Dann gilt:
\[
  3n_l + 1 = 6 \cdot 3^l \cdot m - 2 = 2(3^{l+1} m - 1).
\]
Also $\vii(3n_l + 1) = 1 + \vii(3^{l+1}m - 1)$.
Nach LTE ist $\vii(3^{l+1}m - 1) = O(\log l)$ für typisches~$m$
--- deutlich kleiner als der naive Erwartungswert von~$2$.
\end{tcolorbox}

Die korrekte Analyse muss daher nach $l$ konditionieren.

% ====================================================
\section{Schritt 1: Wahrscheinlichkeiten der Kettenlänge}
\label{sec:wahrscheinlichkeiten}
% ====================================================

\begin{tcolorbox}[ergebnis, title={Bewiesenes Ergebnis -- Kettenlängenverteilung}]
Für eine ungerade Zahl~$n_0$ gilt:
\begin{itemize}[noitemsep]
  \item $l = 0$ genau dann, wenn $n_0 \not\equiv 3 \pmod{4}$
        (kein C-Vorläufer-Typ), d.\,h.\ $P(l=0) = 1/2$.
  \item $l \geq k$ genau dann, wenn $2^{k+1} \mid (n_0+1)$,
        d.\,h.\ $P(l \geq k) = 2^{-k}$.
  \item $P(l = k) = P(l \geq k) - P(l \geq k{+}1) = 2^{-k} - 2^{-(k+1)} = 2^{-(k+1)}$
        für $k \geq 1$.
\end{itemize}
\end{tcolorbox}

\begin{proposition}[Momente der Kettenlänge]
\label{prop:momente}
Unter der obigen geometrischen Verteilung gilt:
\[
  P(l=0) = \tfrac{1}{2}, \qquad
  P(l=k) = 2^{-(k+1)} \text{ für } k \geq 1.
\]
Der Erwartungswert ist
\[
  \E[l] = \sum_{k=1}^{\infty} k \cdot 2^{-(k+1)}
         = \frac{1}{2} \sum_{k=1}^{\infty} k \cdot 2^{-k}
         = \frac{1}{2} \cdot 2 = 1.
\]
\end{proposition}

% ====================================================
\section{Schritt 2: LTE-Schranke für \texorpdfstring{$\nu$}{nu} an
  C-Ketten-Endpunkten}
\label{sec:lte}
% ====================================================

\begin{lemma}[LTE-Schranke, bewiesen]
\label{lem:lte}
Sei $l \geq 1$ und $n_l = 2 \cdot 3^l \cdot m - 1$ mit ungeradem~$m$.
Dann gilt:
\[
  \nu = \vii(3n_l + 1) = 1 + \vii(3^{l+1} m - 1).
\]
Nach dem Lifting-the-Exponent-Lemma für $p=2$ gilt für $2 \nmid m$:
\[
  \vii(3^{l+1} m - 1) \leq \vii(3m-1) + \vii(l+1).
\]
Insbesondere ist $\vii(3^{l+1}m-1) = O(\log l)$ für festes~$m$ und
wachsendes~$l$.
\end{lemma}

\begin{bemerkung}
Konkret bedeutet dies: Für $l \geq 1$ ist der typische Wert von $\nu$
\[
  \E[\nu \mid l=k] \approx 2 + \vii(k+1),
\]
da $\vii(3^{k+1}m-1) \approx 1 + \vii(k+1)$ für generisches~$m$.
Für $k+1$ ungerade gilt $\vii(k+1) = 0$, also $\E[\nu \mid l=k] \approx 2$.
Für $k+1$ gerade gilt $\E[\nu \mid l=k] \approx 2 + \vii(k+1) \geq 3$.
In der Grobabschätzung verwenden wir $\E[\nu \mid l=k] \approx 2$ für alle $k$.
\end{bemerkung}

% ====================================================
\section{Schritt 3: Konditionierter Log-Drift}
\label{sec:konditionierter_drift}
% ====================================================

Der logarithmische Block-Drift ist
\[
  \logF = \log F(B) = (l+1)\log 3 - (l+\nu)\log 2.
\]
Wir berechnen $\E[\logF \mid l=k]$ für $k = 0$ und $k \geq 1$ separat.

\subsection{Fall \texorpdfstring{$l=0$}{l=0}: Direkter EA-Schritt}

Für $l=0$ liegt kein C-Vorläufer vor; $\nu$ ist geometrisch verteilt
mit $P(\nu=j) = 2^{-(j-1)}$ für $j \geq 2$ (mindestens eine freie
Zweierpotenz), also $\E[\nu \mid l{=}0] = 3$.

\begin{proposition}[Drift für $l=0$, heuristisch]
\label{prop:drift_l0}
\[
  \E[\logF \mid l=0]
  = \log 3 - (0+3)\log 2
  = \log_2 3 - 3
  \approx 1{,}585 - 3 = -1{,}415.
\]
\end{proposition}

\subsection{Fall \texorpdfstring{$l=k \geq 1$}{l=k>=1}: Nach C-Kette}

Nach einer C-Kette der Länge $k$ gilt mit $\E[\nu \mid l{=}k] \approx 2$:

\begin{proposition}[Drift für $l=k$, heuristisch]
\label{prop:drift_lk}
\[
  \E[\logF \mid l=k]
  \approx (k+1)\log 3 - (k+2)\log 2
  = k(\log 3 - \log 2) + (\log 3 - 2\log 2).
\]
Mit $\log_2 3 \approx 1{,}585$ und allen Logarithmen zur Basis 2:
\[
  \E[\logF \mid l=k]
  \approx k \cdot 0{,}585 + (1{,}585 - 2)
  = 0{,}585\,k - 0{,}415.
\]
\end{proposition}

\begin{bemerkung}
Man beachte den Vorzeichenwechsel:
\begin{itemize}[noitemsep]
  \item $k=0$: $-0{,}415$ (kontrahierend),
  \item $k=1$: $\phantom{-}0{,}170$ (leicht expandierend!),
  \item $k=2$: $\phantom{-}0{,}755$ (expandierend).
\end{itemize}
Blöcke mit langen C-Ketten expandieren im Mittel.
Der Gesamtdrift hängt davon ab, wie stark die geometrisch abnehmenden
Wahrscheinlichkeiten $P(l=k) = 2^{-(k+1)}$ diese Expansion dämpfen.
\end{bemerkung}

% ====================================================
\section{Schritt 4: Gesamtdrift}
\label{sec:gesamtdrift}
% ====================================================

\begin{proposition}[Gesamtdrift, heuristisch]
\label{prop:gesamtdrift}
\[
  \E[\logF]
  = P(l=0)\cdot\E[\logF \mid l=0]
    + \sum_{k=1}^{\infty} P(l=k)\cdot\E[\logF \mid l=k].
\]
Einsetzen ergibt:
\begin{align*}
  \E[\logF]
  &= \frac{1}{2}\cdot(-1{,}415)
     + \sum_{k=1}^{\infty} \frac{1}{2^{k+1}}\cdot(0{,}585\,k - 0{,}415)\\[4pt]
  &= -0{,}708
     + \frac{1}{2}\Bigl(
         0{,}585\sum_{k=1}^{\infty} k\cdot 2^{-k}
         - 0{,}415\sum_{k=1}^{\infty} 2^{-k}
       \Bigr).
\end{align*}
Mit den Standardwerten $\sum_{k=1}^{\infty} k \cdot 2^{-k} = 2$ und
$\sum_{k=1}^{\infty} 2^{-k} = 1$:
\[
  \E[\logF]
  = -0{,}708 + \frac{1}{2}(0{,}585\cdot 2 - 0{,}415\cdot 1)
  = -0{,}708 + \frac{1}{2}(1{,}170 - 0{,}415)
  = -0{,}708 + 0{,}378
  = -0{,}330.
\]
\end{proposition}

\begin{tcolorbox}[ergebnis, title={Hauptergebnis der konditionierten Rechnung (heuristisch)}]
Unter der Näherung $\E[\nu \mid l{=}k] \approx 2$ für alle $k \geq 1$
und $\E[\nu \mid l{=}0] = 3$ ergibt die konditionierte Driftrechnung:
\[
  \boxed{\E[\logF] \approx -0{,}330 < 0.}
\]
Der logarithmische Drift ist negativ — auch nachdem die Unterschätzung
von~$\nu$ an C-Ketten-Endpunkten korrigiert wurde.
\end{tcolorbox}

\subsection{Tabellarische Übersicht}

\begin{table}[h!]
\centering
\caption{Konditionierte Driftkomponenten pro Kettenlänge~$l$}
\label{tab:driftkomponenten}
\renewcommand{\arraystretch}{1.35}
\begin{tabular}{
  c
  S[table-format=1.4]
  S[table-format=1.1]
  S[table-format=+1.3]
  S[table-format=+1.4]
}
\toprule
$l$ & {$P(l)$} & {$\E[\nu \mid l]$} & {$\E[\logF \mid l]$} & {Beitrag $P(l)\cdot\E[\logF \mid l]$} \\
\midrule
$0$      & 0.5000  & 3.0 & -1.415 & -0.7075 \\
$1$      & 0.2500  & 2.0 & +0.170 & +0.0425 \\
$2$      & 0.1250  & 2.0 & +0.755 & +0.0944 \\
$3$      & 0.0625  & 2.0 & +1.340 & +0.0838 \\
$4$      & 0.0313  & 2.0 & +1.925 & +0.0602 \\
$5$      & 0.0156  & 2.0 & +2.510 & +0.0392 \\
{$\vdots$} & {$\vdots$} & {$\vdots$} & {$\vdots$} & {$\vdots$} \\
{$\infty$} & {(Summe)} & {---} & {---} & {$-0{,}330$} \\
\bottomrule
\end{tabular}
\end{table}

\begin{bemerkung}
Die Beiträge für $l \geq 2$ wachsen linear in~$l$ (im Zähler), aber die
Wahrscheinlichkeiten fallen exponentiell.
Die Reihe konvergiert, und das Gesamtergebnis ist negativ.
\end{bemerkung}

% ====================================================
\section{Schritt 5: Sensitivitätsanalyse}
\label{sec:sensitivitaet}
% ====================================================

\begin{tcolorbox}[sensitivitaet]
\textbf{Frage:} Wie sensitiv ist $\E[\logF]$ gegenüber der Annahme
$\E[\nu \mid l{=}k] = \alpha$ für $k \geq 1$?
\end{tcolorbox}

\subsection{Allgemeine Formel}

Sei $\alpha \geq 1$ ein freier Parameter mit $\E[\nu \mid l{=}k] = \alpha$
für alle $k \geq 1$.
Dann gilt:
\[
  \E[\logF \mid l=k] = (k+1)\log 3 - (k+\alpha)\log 2
  = k(\log 3 - \log 2) + \log 3 - \alpha\log 2.
\]
Also:
\begin{align*}
  \E[\logF](\alpha)
  &= -0{,}708
     + \frac{1}{2}\Bigl(
         (\log 3 - \log 2)\sum_{k=1}^{\infty} k \cdot 2^{-k}
         + (\log 3 - \alpha\log 2)\sum_{k=1}^{\infty} 2^{-k}
       \Bigr)\\
  &= -0{,}708 + \frac{1}{2}\bigl(0{,}585\cdot 2 + (1{,}585 - \alpha)\cdot 1\bigr)\\
  &= -0{,}708 + \frac{1}{2}(1{,}170 + 1{,}585 - \alpha)\\
  &= -0{,}708 + \frac{1}{2}(2{,}755 - \alpha)\\
  &= -0{,}708 + 1{,}378 - \frac{\alpha}{2}\\
  &= 0{,}670 - \frac{\alpha}{2}.
\end{align*}

\begin{proposition}[Kritischer $\alpha$-Wert]
\label{prop:kritisch}
Das Vorzeichen von $\E[\logF](\alpha)$ wechselt bei
\[
  \alpha^* = 2 \cdot 0{,}670 = 1{,}340.
\]
Für $\alpha < \alpha^* \approx 1{,}34$ wäre der Gesamtdrift positiv.
Für $\alpha > \alpha^*$ ist er negativ.
\end{proposition}

\begin{table}[h!]
\centering
\caption{Sensitivitätsanalyse: $\E[\logF]$ als Funktion von $\alpha = \E[\nu \mid l{\geq}1]$}
\label{tab:sensitivitaet}
\renewcommand{\arraystretch}{1.3}
\begin{tabular}{
  S[table-format=1.1]
  S[table-format=+1.4]
  l
}
\toprule
{$\alpha$} & {$\E[\logF](\alpha)$} & Interpretation \\
\midrule
1.0 & +0.170 & stark expansiv (physikalisch ausgeschlossen: $\nu \geq 1$) \\
1.5 & -0.080 & schwach kontraktiv \\
2.0 & -0.330 & kontraktiv (Basisannahme dieses Dokuments) \\
2.5 & -0.580 & stärker kontraktiv \\
3.0 & -0.830 & naive Schätzung (geometrisch, unkonditioniert) \\
3.5 & -1.080 & sehr stark kontraktiv \\
\midrule
{$1{,}34$} & {$\approx 0{,}000$} & \textbf{kritischer Wert $\alpha^*$} \\
\bottomrule
\end{tabular}
\end{table}

\begin{tcolorbox}[warnung, title={Interpretation der Sensitivitätsanalyse}]
\begin{itemize}[noitemsep]
  \item Der LTE-Mechanismus erzwingt $\nu \geq 1$ (trivial) und typisch
        $\nu \geq 2$ an C-Ketten-Endpunkten.
        Damit gilt $\alpha \geq 2$, was den kritischen Wert $\alpha^* \approx 1{,}34$
        weit unterschreitet.
  \item \textbf{Fazit:} Solange die LTE-Schranke $\E[\nu \mid l{\geq}1] \geq 2$
        gilt (was aus der Struktur der C-Ketten folgt), ist $\E[\logF] \leq -0{,}33$.
        Das Ergebnis ist gegenüber plausiblen Variationen in~$\alpha$ robust.
  \item Unsicherheit: Ob $\E[\nu \mid l{=}k]$ genauer bei~$2$ oder bei
        $2{,}5$ liegt, verändert den Drift von $-0{,}33$ auf $-0{,}58$,
        aber nicht das Vorzeichen.
\end{itemize}
\end{tcolorbox}

% ====================================================
\section{Schritt 6: Verbleibende rigoros offene Lücken}
\label{sec:offene_luecken}
% ====================================================

Selbst wenn $\E[\logF] < 0$ korrekt und das heuristische Argument
überzeugend ist, fehlen für einen vollständigen Beweis der Collatz-Vermutung
mindestens zwei entscheidende Schritte.

\begin{tcolorbox}[lücke, title={Offene Lücke 1: Ergodizitätsargument}]
\textbf{Problem:} Der negative Drift $\E[\logF] < 0$ ist ein
\emph{Erwartungswert} über die stationäre Verteilung der Kettenlängen.
Ein Beweis erfordert, dass auch der \emph{Zeitmittelwert}
\[
  \lim_{N\to\infty} \frac{1}{N}\sum_{i=1}^{N} \log F(B_i) < 0
\]
für (fast) jede Bahn gilt.

\textbf{Benötigt:} Ergodiziätssatz (Birkhoff) für die Collatz-Dynamik,
d.\,h.\ Ergodizität des Maßes~$\mu$ bzgl.\ der Abbildung~$\mathcal{U}$.
Die Existenz und Mischungseigenschaften eines solchen invarianten Maßes
sind für die Collatz-Abbildung \textbf{nicht bewiesen}.

\textbf{Konsequenz:} Ohne Ergodizität kann $\E[\logF] < 0$ nur für
``generische'' (zufällig gewählte) Bahnen gelten, nicht für alle
$n \in \N$.
\end{tcolorbox}

\begin{tcolorbox}[lücke, title={Offene Lücke 2: Maß-Null-Ausnahmen}]
\textbf{Problem:} Selbst wenn der Drift negativ ist, könnten
nicht-typische Bahnen existieren:
\begin{enumerate}[noitemsep]
  \item \emph{Nicht-triviale Zyklen:} Bahnen, die in einem Zyklus
        $\neq \{1,2,4\}$ stecken.
        Zwar zeigt $\E[\logF] < 0$ globale Kontraktion,
        aber ein Zyklus kann lokal konsistent sein.
  \item \emph{Divergente Bahnen:} Bahnen mit $T^k(n) \to \infty$.
        Ein negativer mittlerer Drift schließt Divergenz
        \emph{für typische} $n$ aus, nicht formal für alle.
\end{enumerate}

\textbf{Status:}
Nicht-triviale Zyklen sind für $n < 2^{68}$ numerisch ausgeschlossen
(Oliveira e Silva, 2010); für alle $n$ ist es unbewiesen.
Divergenz ist ebenso unbewiesen.

\textbf{Was wäre nötig:}
Ein quantitatives Lyapunov-Argument, das zeigt, dass der Drift
gleichmäßig negativ (nicht nur im Mittel) ist,
oder ein direkter Zyklusausschluss via Algebra.
\end{tcolorbox}

% ====================================================
\section{Zusammenfassung und Abgrenzung}
\label{sec:zusammenfassung}
% ====================================================

\begin{table}[h!]
\centering
\caption{Statusübersicht der Argumentationsschritte}
\label{tab:status}
\renewcommand{\arraystretch}{1.4}
\begin{tabular}{p{5.5cm} c p{6.5cm}}
\toprule
\textbf{Aussage} & \textbf{Status} & \textbf{Anmerkung} \\
\midrule
$P(l=k) = 2^{-(k+1)}$ für $k \geq 1$
  & \textcolor{bewiesengrün}{\textbf{bewiesen}}
  & Folgt aus $2^{k+1} \mid (n_0+1)$ \\
LTE-Schranke $\vii(3n_l+1) = O(\log l)$
  & \textcolor{bewiesengrün}{\textbf{bewiesen}}
  & Standardresultat (Kummer/LTE) \\
$\E[\nu \mid l{=}0] = 3$
  & \textcolor{bewiesengrün}{\textbf{bewiesen}}
  & Geometrische Verteilung, kein C-Vorläufer \\
$\E[\nu \mid l{=}k] \approx 2$ für $k \geq 1$
  & \textcolor{warnorange}{\textbf{heuristisch}}
  & LTE gibt Schranke; exakter Wert unbekannt \\
$\E[\logF] \approx -0{,}33$
  & \textcolor{warnorange}{\textbf{heuristisch}}
  & Folgt aus obiger Näherung \\
Negativität für alle Bahnen
  & \textcolor{lückenrot}{\textbf{offen}}
  & Ergodizität fehlt \\
Zyklenfreiheit / Nicht-Divergenz
  & \textcolor{lückenrot}{\textbf{offen}}
  & Maß-Null-Ausnahmen nicht ausgeschlossen \\
\bottomrule
\end{tabular}
\end{table}

\begin{tcolorbox}[hauptsatz, title={Präzisiertes Hauptresultat (heuristisch)}]
Unter der plausiblen Näherung $\E[\nu \mid l{=}k] \approx 2$ für $k \geq 1$
und $\E[\nu \mid l{=}0] = 3$ gilt für den konditionierten logarithmischen
Block-Drift:
\[
  \E[\logF]
  \;=\; \frac{1}{2}\cdot(-1{,}415)
        + \frac{1}{2}\Bigl(0{,}585\cdot\E[l\,|\,l\geq 1] - 0{,}415\Bigr)
  \;\approx\; -0{,}330.
\]
Die Sensitivitätsanalyse zeigt: Das Ergebnis $\E[\logF] < 0$ ist robust
für alle $\alpha = \E[\nu \mid l{\geq}1] \geq \alpha^* \approx 1{,}34$,
was durch die LTE-Schranke $\alpha \geq 2$ garantiert ist.

\medskip
\textbf{Was dies \emph{nicht} beweist:}
Die Collatz-Vermutung. Fehlende Schritte: Ergodizität (Birkhoff-Theorem
für Collatz), Ausschluss nicht-trivialer Zyklen, Ausschluss divergenter
Bahnen.
\end{tcolorbox}

% ====================================================
\appendix
\section{Rechnung der geometrischen Reihen}
\label{app:reihen}
% ====================================================

Für vollständige Nachvollziehbarkeit: Die verwendeten Reihenformeln.

\begin{align*}
  \sum_{k=1}^{\infty} 2^{-k} &= 1, \\
  \sum_{k=1}^{\infty} k \cdot 2^{-k}
    &= \frac{2^{-1}}{(1-2^{-1})^2} = \frac{1/2}{1/4} = 2, \\
  \sum_{k=1}^{\infty} k^2 \cdot 2^{-k} &= 6
  \quad \text{(für die Varianzabschätzung).}
\end{align*}

Diese Werte sind standardmäßig bekannt; die Rechnungen in den
Propositionen~\ref{prop:momente}, \ref{prop:gesamtdrift}
und~\ref{prop:kritisch} nutzen ausschließlich die ersten beiden.

% ====================================================
\section{Verbindung zu \textit{collatz\_lyapunov.tex}}
\label{app:verbindung}
% ====================================================

Dieses Dokument ist als \emph{Präzisierungs-Appendix} zu
\textit{collatz\_lyapunov.tex} gedacht.
Es kann entweder:
\begin{enumerate}[noitemsep]
  \item als eigenständiges Dokument (Appendix) referenziert werden,
  \item oder als Korrektur/Erweiterung von Abschnitt~3 in
        \textit{collatz\_lyapunov.tex} eingearbeitet werden.
\end{enumerate}

Die wesentliche Botschaft: Der negative Drift $\E[\logF] < 0$
überlebt die Präzisierung durch die konditionierte $\nu$-Verteilung,
aber die Lücken zu einem vollständigen Beweis bleiben dieselben wie
im Hauptdokument.

\end{document}
